\newpage{\ }
\thispagestyle{empty}
\chapter*{Resumen}

\noindent Las pequeñas y medianas empresas (pymes) necesitan cada vez más
aplicaciones web propias, pero rara vez pueden asumir el coste de construir
desde cero los componentes que casi cualquier aplicación SaaS
(\textit{Software as a Service}) repite: autenticación segura, gestión de
usuarios y equipos con roles, cobro de suscripciones y una herramienta para
administrar los datos sin programar. Este Trabajo Fin de Grado diseña,
implementa y valida PymeKit, una arquitectura software reutilizable que reúne
esos componentes en una base común sobre la que desarrollar de forma ágil
nuevas aplicaciones SaaS orientadas a pymes.

% PENDIENTE (P-06): forma de presentar la parte reutilizada, a confirmar con el tutor.
PymeKit no parte de cero. Integra dos bases de código existentes con
licencia ---una base de aplicación SaaS y un gestor de contenidos (CMS) para
la base de datos--- y las adapta y extiende hasta formar un único producto. Las
aportaciones del trabajo son la integración del CMS dentro de la consola de
administración de la aplicación, con un solo servicio, un solo inicio de
sesión y una sola pila de librerías; la reorganización de la consola por
áreas de negocio comprensibles para una pyme; la interfaz íntegramente en
español sobre una infraestructura de internacionalización extensible; la
centralización de la configuración reutilizable; la corrección documentada de
numerosos fallos de seguridad heredados; y un proceso de desarrollo asistido
por agentes de inteligencia artificial, controlado mediante directrices,
comprobaciones automáticas y revisiones adversariales.

La solución se organiza como un monorepo de paquetes desacoplados con
TypeScript estricto. La aplicación web usa TanStack Start y React; los datos,
la autenticación y el almacenamiento se apoyan en Supabase (PostgreSQL); los
pagos, en Stripe; y la API del CMS, en Hono y Drizzle, montada dentro del
propio servidor web. La autorización se aplica en la base de datos mediante
seguridad a nivel de fila (RLS), de modo que el aislamiento entre clientes
(\textit{multi-tenancy}) no depende solo del código de la aplicación. El
acceso a la consola exige un segundo factor de autenticación, y las acciones
realizadas desde el CMS quedan registradas en un registro de auditoría que el
propio personal no puede falsificar.

La calidad se verifica con pruebas unitarias, más de dos mil pruebas de base
de datos con pgTAP ---incluidas pruebas de ataque entre clientes--- y pruebas
de extremo a extremo con Playwright, ejecutadas en integración continua. El
proceso dejó documentadas en una bitácora más de cincuenta incidencias, buena
parte de ellas fallos de seguridad del código de partida que se detectaron y
corrigieron durante la integración. La memoria describe la planificación, el
análisis, el diseño, la construcción, las pruebas y el despliegue de la
plataforma, la evaluación de su reutilización y los manuales de usuario y de
desarrollador.

\vspace{.5cm}

% Añadir una serie de palabras clave que aparezcan en el texto
\noindent \textbf{Palabras clave:} SaaS, pymes, arquitectura reutilizable,
\textit{multi-tenancy}, seguridad a nivel de fila (RLS), Supabase, PostgreSQL,
CMS, TanStack Start, desarrollo asistido por IA.
